\documentclass[english, parskip=full]{scrartcl}
\usepackage[utf8]{inputenc}  % Kodierung der Datei
\usepackage[english]{babel}  % Sprache auf Deutsch 
\usepackage{color}
\usepackage{graphicx, gensymb}
\usepackage{amsfonts, cancel, amssymb}
\usepackage{amsmath, amsthm, array, esint}
\usepackage{booktabs, multirow}  % schönere Tabellen
\usepackage{caption}  % Anpassung der Captions
\usepackage{scrlayer-scrpage}    % Manipulation des Seitenstils
\pagestyle{scrheadings}  % Stil für die Seite setzen
\automark{section}  % Abschnittsnamen als Seitenbeschriftung 
\setheadsepline{.5pt}    % Kopzeile durch Linie abtrennen
\usepackage[hidelinks]{hyperref}  % Links und weitere PDF-Features
\usepackage{typearea, tikz, pgffor, verbatim}
\usetikzlibrary{calc, quotes}
\usepackage[cache=false, outputdir=build]{minted}
\usepackage{placeins} % provides \FloatBarrier

\definecolor{bg}{rgb}{0.9,0.9,0.9}
\newcommand\numberthis{\addtocounter{equation}{1}\tag{\theequation}}
\newcommand{\bra}{}
\newcommand{\kom}{}
\newcommand{\brak}{}
\newcommand{\braket}{}
\newcommand{\intx}{}
\newcommand{\intr}{}
\newcommand{\kla}{}
\newcommand{\klam}{}
\newcommand{\klammer}{}
\newcommand{\oper}{}
\newcommand{\tecxt}{}
\newcommand{\Bra}{}
\newcommand{\Ket}{}
\renewcommand{\i}{\text{i}}
\newcommand{\e}{\text{e}}
\newcommand*{\Fourier}{\textsc{Fourier}\ }
\newcommand*{\F}{\mathcal{F}}
\newcommand*{\sinc}{\text{sinc\,}}
\newcommand{\conv}{\mathop{\scalebox{3.5}{\raisebox{-0.38ex}{\makebox[0.5\width][c]{$\ast$}}}}\limits}

\newcommand*{\titel}{03 Ornstein-Uhlenbeck process}
\newcommand*{\autor}{Joachim Schwardt and Julian Fleck}

\hypersetup{pdfauthor={\autor}, pdftitle={\titel}}

%\titlehead{titlehead}
\subject{Stochastic processes}
\title{\titel}
\author{\autor}
\date{\today}

\begin{document}

%\begin{titlepage}
\maketitle  % Titel setzen
%\tableofcontents  % Inhaltsverzeichnis setzen
%\end{titlepage}

\section*{Problem 3.1: Gaussian Markov process}
We are given a gaussian stochastic process $Y(t)$ with autocorrelation function $\kappa(t_1,t_2)$ and variance $\sigma^2(t)$.
Define the quantity
\begin{align}
\rho(t_1,t_2) & := \frac{\kappa(t_1,t_2)}{\sigma(t_1)\sigma(t_2)}.
\end{align}
Then the condition
\begin{align}
\rho(t_3,t_1) &= \rho(t_3,t_2)\rho(t_2,t_1)
\end{align}
is equivalent to
\begin{align}
\sigma^2(t_2)\kappa(t_3,t_1) &= \kappa(t_3,t_2)\kappa(t_2,t_1).
\end{align}
Recall the definition of the autocorrelation function (assuming $\bra\langle Y(t) \rangle\equiv 0$ for simplicity)
\begin{align}
\kappa (t_1,t_2) &:=\bra\langle Y(t)Y(s) \rangle = \intr\int_{\mathbb{R}^{2}} y_1y_2p_2(y_2,t_2,y_1,t_1) \, \mathrm{d}^{2}y = \\
&=\intr\int_{\mathbb{R}^{2}} y_1y_2p_{1|1}(y_2,t_2|y_1,t_1)p_1(y_1,t_1) \, \mathrm{d}^{2}y.
\end{align}
Since $Y(t)$ is gaussian, the probability distributions are given by
\begin{align}
p_1(y,t) &= \frac{1}{\sqrt{2\pi\sigma^2(t)}}\exp\kla\left( -\frac{y^2}{2\sigma^2(t)} \right),\\
p_2(y_2,t_2,y_1,t_1) &= \frac{1}{2\pi\sqrt{\alpha_{21}}}\exp \kla\left( -\frac{\sigma_1^2\sigma_2^2}{2\alpha_{21}}\klam\left[ \frac{y_2^2}{\sigma_2^2} + \frac{y_1^2}{\sigma_1^2} - 2\frac{\kappa_{21}}{\sigma_1^2\sigma_2^2}y_1y_2 \right] \right),
\end{align}
where $\sigma_j := \sigma(t_j)$, $\kappa_{jk} :=\kappa(t_j,t_k)$ and $\alpha_{jk} := \sigma_j^2\sigma_k^2 - \kappa_{jk}^2$.
Note that the relation $1 - \frac{\alpha_{21}}{\sigma_1^2\sigma_2^2} = \frac{\kappa_{21}^2}{\sigma_1^2\sigma_2^2}$ by construction.
Then the transition probability distribution is
\begin{align}
p_{1|1}(y_2,t_2|y_1,t_1) &= \frac{p_2(y_2,t_2,y_1,t_1)}{p_1(y_1,t_1)} = \\
&= \frac{\sigma_1}{\sqrt{2\pi\alpha_{21}}} \exp\kla\left( -\frac{1}{2\alpha_{21}}\klam\left[ \sigma_1^2y_2^2 + \sigma_2^2y_1^2 - 2\kappa_{21}y_2y_1 - \frac{\alpha_{21}}{\sigma_1^2}y_1^2 \right] \right) = \\
&= \frac{\sigma_1}{\sqrt{2\pi\alpha_{21}}} \exp\kla\left( -\frac{1}{2\alpha_{21}}\klam\left[ \sigma_1^2y_2^2 - 2\kappa_{21}y_2y_1 + \sigma_2^2y_1^2\frac{\kappa_{21}^2}{\sigma_1^2\sigma_2^2} \right] \right) =  \\
&= \frac{\sigma_1}{\sqrt{2\pi\alpha_{21}}} \exp\kla\left( -\frac{1}{2\alpha_{21}}\klam\left[ \sigma_1y_2 - \frac{\kappa_{21}}{\sigma_1}y_1 \right]^2 \right).
\end{align}
We can now check the Chapman-Kolmogorov equation, i.e.
\begin{align*}
&\intr\int_{\mathbb{R}^{ }} p_{1|1}(y_3,t_3|y_2,t_2)p_{1|1}(y_2,t_2|y_1,t_1) \, \mathrm{d}^{ }y_2 =  \\
&= \frac{\sigma_1\sigma_2}{2\pi\sqrt{\alpha_{32}\alpha_{21}}} \intr\int_{\mathbb{R}^{ }} \exp\kla\left( -\frac{1}{2\alpha_{32}}\klam\left[ \sigma_2y_3 - \frac{\kappa_{32}}{\sigma_2}y_2 \right]^2 -\frac{1}{2\alpha_{21}}\klam\left[ \sigma_1y_2 - \frac{\kappa_{21}}{\sigma_1}y_1 \right]^2 \right) \, \mathrm{d}^{ }y_2 =  \numberthis \\
&= \frac{\sigma_1\sigma_2}{2\pi\sqrt{\alpha_{32}\alpha_{21}}} \exp\kla\left( -\frac{\sigma_2^2y_3^2}{2\alpha_{32}}-\frac{\kappa_{21}^2y_1^2}{2\alpha_{21}\sigma_1^2} \right) \cdot\\
&\hspace{1cm} \cdot \intr\int_{\mathbb{R}^{ }} \exp\kla\left( -\frac{y_2^2}{2}\klam\left[ \frac{\kappa_{32}^2}{\sigma_2^2\alpha_{32}} + \frac{\sigma_1^2}{\alpha_{21}} \right] + y_2\klam\left[ \frac{\kappa_{32}}{\alpha_{32}}y_3 + \frac{\kappa_{21}}{\alpha_{21}}y_1 \right] \right) \, \mathrm{d}^{ }y_2 =  \numberthis \\
&= \frac{\sigma_1\sigma_2}{2\pi\sqrt{\alpha_{32}\alpha_{21}}}\sqrt{\frac{2\pi}{\frac{\kappa_{32}^2}{\sigma_2^2\alpha_{32}} + \frac{\sigma_1^2}{\alpha_{21}} }}\exp\kla\left( -\frac{\sigma_2^2y_3^2}{2\alpha_{32}}-\frac{\kappa_{21}^2y_1^2}{2\alpha_{21}\sigma_1^2} \right) \exp\kla\left( \frac{\left[ \frac{\kappa_{32}}{\alpha_{32}}y_3 + \frac{\kappa_{21}}{\alpha_{21}}y_1  \right]^2}{2\left[ \frac{\kappa_{32}^2}{\sigma_2^2\alpha_{32}} + \frac{\sigma_1^2}{\alpha_{21}} \right]} \right) = \numberthis\\
&= \frac{\sigma_1\sigma_2^2}{\sqrt{2\pi\klam\left[ \alpha_{21}\kappa_{32}^2 + \sigma_1^2\sigma_2^2\alpha_{32} \right]} } \exp\kla\left( -\frac{1}{2}\klam\left[ \frac{\sigma_2^2y_3^2}{\alpha_{32}}+ \dots \right] \right)  \overset{!}{=} \numberthis  \\
&\overset{!}{=} p_{1|1}(y_3,t_3|y_1,t_1) = \frac{\sigma_1}{\sqrt{2\pi\alpha_{31}}} \exp\kla\left( -\frac{1}{2\alpha_{31}}\klam\left[ \sigma_1y_3 - \frac{\kappa_{31}}{\sigma_1}y_1 \right]^2 \right). \numberthis
\end{align*}
Expanding the exponent is a task better suited for algebra programs.
However, we may as well trust that the equation is fulfilled if and only if the prefactor matches, since it is essentially independent of the exponential factor in this case.
Thus we arrive at the condition
\begin{align}
\sigma_2^4\alpha_{31} &\overset{!}{=} \alpha_{21}\kappa_{32}^2 + \sigma_1^2\sigma_2^2\alpha_{32}\\
\Rightarrow\ \ \ \sigma_2^4\kla\left( \sigma_1^2\sigma_3^2 - \kappa_{31}^2 \right) &= \kla\left( \sigma_1^2\sigma_2^2 - \kappa_{21}^2 \right)\kappa_{32}^2 + \sigma_1^2\sigma_2^2\kla\left( \sigma_2^2\sigma_3^2 - \kappa_{32}^2 \right) \\
\Rightarrow\ \ \ -\kappa_{31}^2\sigma_2^4 &= -\kappa_{21}^2\kappa_{32}^2.
\end{align}
This is of course equivalent to
\begin{align}
\sigma_2^2 \kappa_{31} &= \kappa_{32}\kappa_{21} \label{eqn:condition gauss markov} ,
\end{align}
thus this condition is indeed necessary for $Y(t)$ to be a Markov process, since it is necessary to fulfil the Chapman-Kolmogorov equation.
However, for it to also be sufficient, we now have to compare the exponents neglected above.
Since we can already assume the identity, this is now a little bit less cumbersome.
Comparing the exponents multiplied by $-2$ for simplicity gives
\begin{align}
\tecxt\text{-2e} &:= \frac{\sigma_2^2y_3^2}{\alpha_{32}} + \frac{\kappa_{21}^2y_1^2}{\alpha_{21}\sigma_1^2} - \frac{\kla\left( \alpha_{21}\kappa_{32}y_3 + \alpha_{32}\kappa_{21}y_1 \right)^2}{\kla\left( \frac{\alpha_{21}\kappa_{32}^2}{\sigma_2^2} + \sigma_1^2\alpha_{32} \right)\alpha_{32}\alpha_{21}} = \\
&= y_3^2\kla\left( \frac{\sigma_2^2}{\alpha_{32}} - \frac{\alpha_{21}\kappa_{32}^2}{\sigma_2^2\alpha_{31}\alpha_{32}} \right) + y_1^2\kla\left( \frac{\kappa_{21}^2}{\alpha_{21}\sigma_1^2} - \frac{\alpha_{32}\kappa_{21}^2}{\sigma_2^2\alpha_{31}\alpha_{21}} \right) - 2y_1y_3\frac{ \kappa_{32}\kappa_{21} }{\sigma_2^2\alpha_{31}} = \\
&= \frac{1}{\alpha_{31}}\kla\left( \frac{y_3^2}{\alpha_{32}\sigma_2^2}\klam\left[ \sigma_2^4\alpha_{31} - \alpha_{21}\kappa_{32}^2 \right] + \frac{y_1^2\kappa_{21}^2}{\alpha_{21}\sigma_1^2}\klam\left[ \alpha_{31} - \frac{\sigma_1^2\alpha_{32}}{\sigma_2^2} \right] - 2y_1y_3\kappa_{31} \right) = \\
&= \frac{1}{\alpha_{31}}\kla\left( \sigma_1^2y_3^2 + \frac{y_1^2\kappa_{21}^2\kappa_{32}^2}{\sigma_1^2\sigma_2^4} - 2y_1y_3\kappa_{31} \right) = \\
&= \frac{1}{\alpha_{31}}\kla\left( \sigma_1^2y_3^2 + \frac{y_1^2\kappa_{31}^2}{\sigma_1^2} - 2y_1y_3\kappa_{31} \right) = \\
&= \frac{1}{\alpha_{31}}\kla\left( \sigma_1y_3-\frac{\kappa_{31}}{\sigma_1}y_1 \right)^2,
\end{align}
which is exactly the expected result for the exponent of $p_{1|1}(3|1)$.
Therefore \eqref{eqn:condition gauss markov} is indeed also sufficient for $Y(t)$ to be a Markov process.\\
For the Wiener process we have $\sigma^2(t) = t$ and $\kappa (t,s) = \min\{t,s\}$.
Since \eqref{eqn:condition gauss markov} only contains autocorrelation functions with time-ordered arguments, we simply find
\begin{align}
\kappa(t_3,t_2)\kappa(t_2,t_1) &= t_2t_1 = \sigma^2(t_2)\kappa(t_3,t_1).
\end{align}  
Thus the Wiener process obeys this relation, which was to be expected.
\newpage


\section*{Problem 3.2: Ornstein-Uhlenbeck process I}
Let $Y(t)$ be the standard Ornstein-Uhlenbeck process, i.e.
the (unique) stationary, gaussian Markov process with autocorrelation function
\begin{align}
\kappa(\tau) &= \e^{-\tau}
\end{align}
for $\tau = t_2 - t_1 > 0$ and variance $\sigma^2(t)\equiv 1$. 
Define the integral of this process as
\begin{align}
Z(t) &:= \intx\int_{0}^{t}Y(s)\, \mathrm{d}s.
\end{align}
We find the autocorrelation of $Z(t)$ by first considering $s\ge t$
\begin{align}
\kappa_Z(t,s) &:= \bra\langle Z(t)Z(s) \rangle = \intx\int_{0}^{t}\intx\int_{0}^{s}\bra\langle Y(t')Y(s') \rangle\, \mathrm{d}s'\, \mathrm{d}t' = \\
&= \intx\int_{0}^{t}\intx\int_{0}^{s}\e^{-|t'-s'|}\, \mathrm{d}s'\, \mathrm{d}t' = \\
&=\intx\int_{0}^{t}\kla\left( \intx\int_{0}^{t'}\e^{-(t'-s')}\, \mathrm{d}s' + \intx\int_{t'}^{s}\e^{-(s'-t')}\, \mathrm{d}s' \right)\, \mathrm{d}t' = \\
&= \intx\int_{0}^{t} 1 - \e^{-t'} - \e^{t'-s} + 1\, \mathrm{d}t' = \\
&= 2t + \e^{-t} - 1 - \e^{t-s} + \e^{-s}.
\end{align}
Now assuming $s\ge t$ is equivalent to substituting $s\rightarrow \max\{s,t\}$ and $t\rightarrow\min\{s,t\}$ in the final result, giving
\begin{align}
\kappa_Z(t,s) &= \e^{-t} + \e^{-s} - 1 - \e^{-|t-s|} + 2\min\{t,s\}.
\end{align}
This formula is now symmetric in $t$ and $s$, as expected.
Note that the process is not stationary, since $\kappa_Z(t,s)$ can not be written as a function of the difference $\tau := t-s$.
However, $Z(t)$ is gaussian, since it is an integral of a gaussian process.
%(Finite case is clear, does this really hold for ``continuous sums'' in general, and how to prove?)
The variance can be calculated in a similar manner, since
\begin{align}
\sigma_Z^2(t) &= \bra\langle Z(t)^2 \rangle = \intx\int_{0}^{t}\intx\int_{0}^{t}\bra\langle Y(t_1)Y(t_2) \rangle\, \mathrm{d}t_1\, \mathrm{d}t_2 = \\
&= \intx\int_{0}^{t}\intx\int_{0}^{t}\e^{-|t_2-t_1|}\, \mathrm{d}t_1\, \mathrm{d}t_2 = \\
&= 2\intx\int_{0}^{t}\intx\int_{0}^{t_2}\e^{-(t_2-t_1)} \, \mathrm{d}t_1\, \mathrm{d}t_2 = \\
&= 2\intx\int_{0}^{t} 1 - \e^{-t_2}\, \mathrm{d}t_2 = 2t + 2\e^{-t} - 2.
\end{align}
Here we made use of the symmetry along the line $t_1=t_2$, allowing us to only compute the integral over a triangle ($t_2\in [0,1]$ and $0\le t_1 \le t_2$).
(NOTE: Of course this could just as well be obtained from $\kappa_Z(t,t)$, but now I have already written down the calculation...)
\\
We can use the condition from problem 3.1 to check, if the process is markovian.
Expanding the terms, it quickly becomes clear, that the equality required by \eqref{eqn:condition gauss markov} is not obeyed.
\begin{align*}
\kappa_Z(t_3,t_2)\kappa_Z(t_2,t_1) &= \kla\left( \e^{-t_3} + \e^{-t_2} - 1 - \e^{t_2 - t_3} + 2t_2 \right) \cdot \\
&\hspace{0.7cm} \cdot \kla\left( \e^{-t_2} + \e^{-t_1} - 1 - \e^{t_1 - t_2} + 2t_1 \right) = \numberthis \\
&= \e^{-t_3-t_2} + \e^{-t_3-t_1} - \e^{-t_3} - \e^{-t_3+t_1-t_2} + 2t_1\e^{-t_3} + \\
&\hspace{0.7cm} + \e^{-2t_2} + \e^{-t_2-t_1} - \e^{-t_2} - \e^{-2t_2+t_1} + 2t_1\e^{-t_2} - \\
&\hspace{0.7cm} - \e^{-t_2} - \e^{-t_1} {\color{blue}{+ 1}} + \e^{t_1 - t_2} - 2t_1 - \\
&\hspace{0.7cm} - \e^{-t_3} + \e^{t_2-t_3-t_1} - \e^{t_2-t_3} - \e^{t_1 - t_3} + 2t_1\e^{t_2-t_3} + \\
&\hspace{0.7cm} + 2t_2\e^{-t_2} + 2t_2\e^{-t_1} - 2t_2 - 2t_2\e^{t_1 - t_2} + 4t_2t_1 = ,\numberthis\\
\sigma_Z^2(t_2)\kappa_Z(t_3,t_1) &= 2\kla\left( t_2 - 1 + \e^{-t_2} \right)\kla\left( \e^{-t_3} + \e^{-t_1} - 1 - \e^{t_1 - t_3} + 2t_1 \right) = \numberthis \\
&= 2\e^{-t_3-t_2} + 2\e^{-t_1-t_2} - 2\e^{-t_2} - 2\e^{t_1-t_3-t_2} + 4t_1\e^{-t_2} + \\
&\hspace{0.7cm} - 2\e^{-t_3} - 2\e^{-t_1} {\color{blue}{+ 2}} + 2\e^{t_1 - t_3} - 4t_1 + \\
&\hspace{0.7cm} + 2t_2\e^{-t_3} + 2t_2\e^{-t_1} - 2t_2 - 2t_2\e^{t_1 - t_3} + 4t_2t_1 \numberthis.
\end{align*}
Some terms are equal, but the very simple comparison of the highlighted terms ${\color{blue}{+2}}$ against ${\color{blue}{+ 1}}$ shows that not everything cancels out.
Therefore $Z(t)$ is not a Markov process.

\newpage
\section*{Problem 3.3: Ornstein-Uhlenbeck process II}
Since $Z(t)$ is gaussian, only the first two cumulants contribute to the characteristic functional.  
Thus it is given by 
\begin{align}
\log G[k(t)] &= \i\intr\int_{\mathbb{R}^{ }} k(t)\underbrace{\bra\langle Z(t) \rangle}_{\equiv\, 0} \, \mathrm{d}^{ }t - \frac{1}{2}\intr\int_{\mathbb{R}^{2}} k(t)k(s)\kappa_Z(t,s) \, \mathrm{d}^{}(t,s) = \\
&= -\intr\int_{\mathbb{R}^2} k(t)k(s)\kla\left( \e^{-t} + \e^{-s} - 1 - \e^{-|t-s|} + 2\min\{t,s\} \right) \, \mathrm{d}^{}(t,s).
\end{align}
The general formulation is
\begin{align}
G_Z[k] &= \left\bra\langle \exp\kla\left( \i\intr\int_{\mathbb{R}^{ }} k(t)Z(t) \, \mathrm{d}^{ }t \right) \right\rangle.
\end{align}
Note that $k(t) := \delta(t-t_1) - \delta(t-t_2)$ gives
\begin{align}
\intr\int_{\mathbb{R}^{ }} \klam\left[ \delta(t-t_1) - \delta(t-t_2) \right]Z(t) \, \mathrm{d}^{ }t = Z(t_1) - Z(t_2), 
\end{align}
which leads to
%\begin{align*}
%\left\langle \cos\kla\left( Z(t_1) - Z(t_2) \right) \right\rangle &= G_Z[\delta(t-t_1) - \delta(t-t_2)] =  \numberthis \\
%&= \exp\kla\bigg( -\intr\int_{\mathbb{R}^2} [\delta(t-t_1) - \delta(t-t_2)] [\delta(s-t_1) - \delta(s-t_2)]\cdot\\
%&\hspace{1.2cm} \cdot \kla\left( \e^{-t} + \e^{-s} - 1 - \e^{-|t-s|} + 2\min\{t,s\} \right) \, \mathrm{d}(t,s) \bigg) = \numberthis \\
%&= \exp\kla\Big( \e^{-t_1} + \e^{-t_2} - 1 - \e^{t-s} + 2t \Big)• \numberthis 
%\end{align*}
\begin{align}
\left\langle \cos\kla\left( Z(t_1) - Z(t_2) \right) \right\rangle &= G_Z[\delta(t-t_1) - \delta(t-t_2)] =   \\
&= \exp\kla\bigg( -\frac{1}{2}\intr\int_{\mathbb{R}^2} k(t)k(s) \kappa_Z(t,s) \, \mathrm{d}(t,s) \bigg) =  \\
&= \exp\kla\Big( -\frac{1}{2}\kappa_Z(t_1,t_1) - \frac{1}{2}\kappa_Z(t_2,t_2) + \kappa_Z(t_1,t_2) \Big) = \\
&= \exp\kla\left( 1 - t_1 - t_2 - \e^{-|t_1-t_2|} + 2\min\{t_1,t_2\}  \right)  = \\
&= \exp\kla\left( 1 - |t_1 - t_2| - \e^{-|t_1-t_2|} \right),
\end{align}
where we used $\min\{2x,2y\}-x-y = \min\{x-y,y-x\} = -|x-y|$.


\newpage


\section*{Problem 3.4: Ornstein-Uhlenbeck process III}
Let $Y(t)$ be the standardized Ornstein-Uhlenbeck process and $t > t_n > \dots > t_1$.
Let $E(t,t_n,\dots,t_1) := \bra\langle Y(t)Y(t_n)\dots Y(t_1) \rangle$, which gives
\begin{align}
\frac{\tecxt\text{d}}{\tecxt\text{d}t}E &= \intr\int_{\mathbb{R}^{n+1}} yy_n\dots y_1 \frac{\tecxt\text{d}}{\tecxt\text{d}t} p_{n+1}(y,t,y_n,t_n,\dots ,y_1,t_1 \, \mathrm{d}^{}(y,y_n, \dots, y_1) = \\
&= \intr\int_{\mathbb{R}^{n }} y_n\dots y_1p_n(y_n,\dots,y_1)\intr\int_{\mathbb{R}^{ }} y\partial_t p_{1|1} (y,t|y_n,t_n) \, \mathrm{d}^{ }y \, \mathrm{d}^{}(y_n,\dots, y_1) \label{eqn:time derivative of expectation value}.
\end{align}
Recall the diffusion equation for the transition probability
\begin{align}
\partial_t p_{1|1} &= \partial_y\kla\left( yp_{1|1} \right) + \partial_y^2 p_{1|1}.
\end{align}
We can insert this into the integral expression above and use integration by parts.
Note that the boundary terms both vanish, assuming that $p_{1|1}$ decays faster than $y^2$, and $\partial_yp_{1|1}$ decays faster than $y$ respectively.
Since $p_{1|1}$ is a gaussian, this is justified.
\begin{align}
I &:= \intr\int_{\mathbb{R}^{ }} y\partial_t p_{1|1} \, \mathrm{d}^{ }y = \intr\int_{\mathbb{R}^{ }} y\partial_y\kla\left( yp_{1|1} \right) + y\partial_y^2 p_{1|1} \, \mathrm{d}^{ }y = \\
&= y^2p_{1|1}\Big\vert_\mathbb{R} - \intr\int_{\mathbb{R}^{ }} yp_{1|1} \, \mathrm{d}^{ }y + y\partial_yp_{1|1} \Big\vert_\mathbb{R} - \intr\int_{\mathbb{R}^{ }} \partial_yp_{1|1} \, \mathrm{d}^{ }y = \\
&= -\intr\int_{\mathbb{R}^{ }} yp_{1|1} \, \mathrm{d}^{ }y.
\end{align}
Inserting this back into \eqref{eqn:time derivative of expectation value} gives
\begin{align}
\frac{\tecxt\text{d}}{\tecxt\text{d}t} E &= \intr\int_{\mathbb{R}^{n}} y_n\dots y_1 p_n(y_n\dots y_1) \kla\left( -\intr\int_{\mathbb{R}^{ }} yp_{1|1}(y,t|y_n,t_n) \, \mathrm{d}^{ }y \right)   \, \mathrm{d}^{ }(y_n\dots y_1) = \\
&= -\intr\int_{\mathbb{R}^{n+1}} yy_n\dots y_1p_{n+1}(\dots) \, \mathrm{d}^{ }(y,y_n,\dots,y_1) = -E.
\end{align}
Now let $\Phi(t,[Y])$ be a functional depending on $t$ and all values of $Y$ at smaller times.
Using the product rule we obtain
\begin{align}
\frac{\tecxt\text{d}}{\tecxt\text{d}t}\left\langle Y(t)\Phi(t,[Y]) \right\rangle &= \frac{\tecxt\text{d}}{\tecxt\text{d}t}\langle \overset{\downarrow}{Y}(t) \Phi(t,[Y]) \rangle + \frac{\tecxt\text{d}}{\tecxt\text{d}t}\langle Y(t)\overset{\downarrow}{\Phi}(t,[Y]) \rangle = \\
&= -\left\langle Y(t)\Phi(t,[Y]) \right\rangle + \left\langle Y(t)\partial_t\Phi(t,[Y]) \right\rangle.
\end{align}

\newpage
\section*{Problem 3.5: Ornstein-Uhlenbeck process IV}
Consider the transition probability density with new coordinates $y\rightarrow \alpha y$ and $t\rightarrow \beta t$
\begin{align}
p_{1|1}(y_2,t_2,y_1,t_1) &= \frac{\alpha}{\sqrt{2\pi\klam\left[ 1 - \e^{-2\beta (t_2-t_1)} \right]}} \exp\kla\left( -\frac{\alpha^2\klam\left[ y_2 - \e^{-\beta(t_2 - t_1)} y_1 \right]^2}{2\klam\left[ 1 - \e^{-2\beta (t_2-t_1)} \right]} \right).
\end{align}
Note that the additional factor of $\alpha$ in the normalization is a result of considering $p_{1|1}\,\mathrm{d}y$, where $\mathrm{d}y\rightarrow\alpha \mathrm{d}y$, and then absorbing this factor into $p_{1|1}$.
Now consider the limit $\beta\rightarrow 0$, which gives
\begin{align}
1 - \e^{-2\beta (t_2-t_1)} &= 2\beta (t_2 - t_1) + \mathcal{O}(\beta^2).
\end{align}
Inserting this into the transition probability yields
\begin{align}
p_{1|1}(y_2,t_2,y_1,t_1) &\sim \frac{\alpha}{\sqrt{4\pi\beta (t_2 - t_1)}} \exp\kla\left( -\frac{\alpha^2\klam\left[ y_2 - (1 - \beta (t_2 - t_1)) y_1 \right]^2}{4\beta (t_2 - t_1)} \right)
\end{align}
as $\beta\rightarrow 0$.
Now letting $\alpha = \sqrt{2\beta}$ we obtain the Wiener process
\begin{align}
p_{1|1}(y_2,t_2,y_1,t_1) &\sim \frac{1}{\sqrt{2\pi (t_2 - t_1)}} \exp\kla\left( -\frac{\klam\left[ y_2 - (1 - \beta (t_2 - t_1)) y_1 \right]^2}{2 (t_2 - t_1)} \right)\\
&\overset{\,\beta\rightarrow \, 0}{\longrightarrow} \ p^\tecxt\text{(W)}_{1|1}(y_2,t_2,y_1,t_1) = \frac{1}{\sqrt{2\pi (t_2 - t_1)}} \exp\kla\left( -\frac{\left( y_2 - y_1 \right)^2}{2 (t_2 - t_1)} \right).
\end{align}
For the general Ornstein-Uhlenbeck process, the variance of the transition probability approaches 1 for large $\Delta t$, where the deviation from 1 is decaying exponentially.
By taking $\beta\rightarrow 0$, we consider sufficiently short time scales, such that the variance grows linearly with $\Delta t$.
(Note: we have a hunch that we are missing something here...)
%Interpretation (?): $\alpha\sim\sqrt{\beta}$ is again reminiscent of $y\sim\sqrt{t}$, which relates to the typical ``Brownian behaviour'' of $\Delta y^2 \sim \Delta t$.
%But this is not really an interpretation, since we expected this result from the construction. 

%\begin{thebibliography}{9}
%\bibitem{numerical levy stable distributions} S.Ament, M. O'Neil: Accurate and efficient numerical calculation of stable
%densities via optimized quadrature and asymptotics, \url{https://arxiv.org/pdf/1607.04247.pdf}, 31.\,August~2021.
%\end{thebibliography}

\end{document}